%%----------Chapter 4------------------------------------------
\chapter{Literature Review}\label{chap:lit}

In this section, the existing work that applies hybrid systems (including but not limited to \textit{hybrid automata}) to the task of modelling cellular metabolism is reviewed and a number of core conclusions are drawn. Many of these approaches incorporate discrete switching to model regulatory mechanisms, as a discrete state changed followed by shift in dynamics, triggered in response to threshold being crossed for some continuous variable, is a natural framework to represent a change in gene expression (and hence, active metabolic pathways) by means of cellular regulation. In fact, as the internal control mechanism of the cell, cellular regulation provides a natural area for the application of control theoretic techniques to better understand intracellular behavior.

In [Modeling Treatment and Drug Effects at The Molecular Level Using
Hybrid System Theory], a gene regulation network is modelled using a system of piecewise linear/affine equations. Switching dynamics are provided by unit step functions embedded in the continuous equations, which in turn regulate the levels of gene expression in response to their activation/deactivation. The dynamics of each gene are enhanced with the addition with an external drug term, which represents the level at which the external substance increases/decreases expression of that gene. Simulations can then be run, analyzing the change in dynamics in response to the addition of the drug term for \textit{in silico} drug discovery. This model, however, does not directly consider metabolism and only includes regulatory mechanisms known at the time of formulating the model for some cell. Hence, it provides little insight into the nature of the metabolic shifts that are induced in response to a change in the extracellular medium.

In [Modeling and simulation of metabolic networks for estimation of
biomass accumulation parameters], steady-state internal metabolism is coupled with external uptake/secretion dynamics and the behavior of the cell is formulated as a two-location state machine. The internal metabolism is defined by the steady-state stoichiometric equation $Sv = 0$ which is then solved with respect to an objective function. The uptakes/secretions are then determined by the internal flux values selected by the objective function and are bounded by kinetic constraints. However, instead of using the two-step iterative process to simulate metabolism as is done in dFBA, the internal and external dynamics are coupled into a differential-algebraic system of equations (DAE) which is directly solved for. This is a fairly generic approach that requires a limited amount of information about the organism to be known ahead of time and in theory could be applied to simulate a wide variety of organisms. However, solving the system is challenging and cannot be solved with standard solvers; the authors had to use a modified form that uses approximations of the discontinuous terms to solve it in MATLAB. Hence, this approach does not seem to scale well to arbitrary systems as there are no guarantees that the approximations that will need to be made to larger systems with a greater number of discontinuities will faithfully represent the behavior of the original system in question. Furthermore, the system does not provide a way to isolate the different regulatory states that may be induced by shifts in the metabolic network.

In [Formalizing metabolic-regulatory networks by hybrid
automata], metabolic and regulatory dynamics are combined to build a hybrid automaton model. Each discrete location is used to represent a specific regulatory mode, wherein the intra-cellular concentration dynamics are given by Michaelis-Mentens kinetic equations. Different regulatory modes represent changes in gene expression by "turning on/off" specific Michaelis-Menten terms in the concentration dynamics and the regulatory shifts are induced by the concentration of some metabolite crossing an established guard, triggering the change from one location to the next. This model is interesting as it uses the standard hybrid automata formalism to faithfully represent metabolic dynamics with a regulation component. Furthermore, the results were able to be duplicated by creating an identical hybrid automaton in the Ariadne library and running a similar set of simulations over it. These are an attractive set of propositions, but the weakness of this modelling approach comes from its inherent lack of generalizability. For a larger system, there is no efficient way to determine proper coefficients for every reaction to plug into the associated kinetic term. Furthermore, the metabolic pathways in the paper are not genome-scale, but rather made with a drastic set of simplifying assumptions to reduce the number of reactions (and hence, kinetic terms) that will need to have coefficients determined. This presents a trade-off between an accurate model that's impossible to practically assemble and an simplified one that may lose granularity, especially when considering more complex systems. Additionally, like other hybrid systems models that incorporate regulation, this approach assumes that regulatory modes are known before assembling the model and offers no ability to unravel their structure as emergent properties of the metabolic network.

In [Symbolic Systems Biology: Hybrid Modeling
and Analysis of Biological Networks], the authors present the notion of "symbolic systems biology", a heuristic for representing and reasoning about metabolic pathways focusing on their qualitative, instead of quantitative, behavior. This is introduced because highly quantitative models are often difficult to appropriately determine in practice, as determining the ``real" kinetic constants for each reaction a daunting if not computationally intractable task. In face of this difficulty, many models in practice determine the kinetic constants by fitting the model as a whole to experimentally gathered data, using min-sum-of-squares or similar model fitting techniques. The quasi-experimentally determined models are often highly sensitive to variance and noise in the determined parameters and there is no inherent guarantee that they faithfully represent the characteristics of the underlying biological system. Hence, the authors determine that it is preferable to have an inherently meaningful model and strive to be qualitatively, instead of quantitatively, correct in terms of its implementation and interpretation. 

They build two examples of this approach, the first of which represents human glucose metabolism with a piece-wise linear system and uses bounded nondeterminism to account for experimental under-determination. This model takes insulin levels as environmental inputs and uses compartmental representation for the different physiological components that use insulin levels as a trigger start a cascade of glucose uptake. The piece-wise linear system is an abstraction of an underlying nonlinear model that contains switching behavior using sigmoidal functions. The authors then perform a reachable set computation to determine bounds for the output of the first segment of the network, that responds to insulin. Using the stability theory of linear differential equations, eigenvalues for the system are computed which show that the dynamics of the second half of the network are stable under both the upper bound and lower computed in the latter segment. This network also provides an \textit{in silico} model for testing insulin injections on a diabetic patient, which can be simulated by setting the biological insulin level to zero and computing the reachable set after supplemental insulin is added to the system.

, and such models are often highly sensitive to variations in

In [Hybrid Modeling and Simulation of Biomolecular Networks], the authors detail a domain-specific modelling language that incorporates hybrid dynamics, nondeterminism, and stochastic variables to model a protein network that self-regulates through a closed feedback loop. Continuous dynamics are applied to model mRNA concentrations when they are above a determined level, wherein individual particle dynamics and collisions can be ignored in favor of treating the mRNA in the medium as a fully-diffused mixture. At lower levels, however, this continuous, deterministic approximation does not faithfully represent the discrete nature of the particles and the unpredictability of their collisions with one another. Hence, when the mRNA concentration drops below the aformentioned level, the continuous dynamics are substituted in favor of a stochastic model of the interaction behavior between protein and mRNA with updates to the discrete concentrations over a fixed time-step. Another example is also presented that models quorum sensing among a bacterial culture. Intra- and extra-cellular continuous dynamics are coupled with discrete switching behavior that is triggered in response to extracellular concentration of a signalling molecule. Hence, the microbes in the culture are able to collectively detect when they are in ``quorum" by the extracellular concentrations which then governs their internal dynamics (they have three disjoint internal states).


Given the broad areas of application, much work has already been performed towards modeling the metabolism of microbial co-cultures using FBA and other constraint-based methods. However, there is a lack of literature pertaining to the use of hybrid systems to model microbial metabolism, much less the specific (and more challenging) case of microbial co-cultures. The majority of the existing work applying hybrid systems (hybrid automata in particular) to the modeling and analysis of microbial metabolism is limited in scope as well as potential application. This is because 

\begin{enumerate}
    \item Most of the approaches to model metabolism using hybrid systems involve directly modeling the detailed enzyme kinetics. For most organisms, the mechanisms governing each reaction are not known, and it is very unlikely that a complete characterization of all enzyme-catalyzed reactions will be found for any sufficiently large number of microbial organisms under consideration. Even for a single organism of interest, with hundreds to thousands of reactions, elucidating and quantifying the enzymatic mechanism for each reaction becomes infeasible and intractable.
    
    \item Due to the latent difficulty of appropriately characterizing the mechanisms underlying each reaction, the network itself can be simplified on the order of phenotypically-distinct pathways. This renders a kinetic model where entire pathways, instead of reactions, are the basic unit for which kinetic equations are defined. While workable for organisms with well-understood metabolic pathways and regulatory mechanisms, the same problems as the reaction-centric approach are invoked upon generalizing to. Namely, the kinetic equations themselves
    
    Hence, simplifying assumptions about mechanisms are usually made which contradicts the purpose of a detailed kinetic model. Two key objections to this type of kinetic modeling are that the transferability of equation forms among organisms is unknown and the even the transferability of enzyme kinetic behavior between the test tube and the cellular interior is debatable.
    % Moreover, there is variation of enyzme kinetics between different genetic polymorphisms and such variation can only be modeled by combining the kinetic metabolic model with gene regulatory model which significantly increases the complexity of the model
    
    
    \item The problem of modeling the detailed enzyme kinetics is further complicated when a co-culture is considered, where in addition to the reactions within each organism, the detailed dynamics of the metabolites in the medium also needed to be tracked. 
    % \textbf{Comment (JFB): we track the dynamics of the concentration within the medium in our approach, and it is also done in dFBA.}
\end{enumerate}

In light of the above observations, it can be concluded that hybrid systems approaches to model microbial metabolism either require either a great deal of preexisting knowledge of reaction mechanisms, drastic simplifying assumptions, or both. Hence, the use of hybrid systems approaches to develop models for metabolism has not seen much popularity. Since the lack of knowledge on enzyme kinetic mechanisms and related parameters can be bypassed using the constrained-based analysis methods that utilize reaction stoichiometries instead of enzyme kinetics, it is worthwhile to develop an approach that builds hybrid systems models for a microbe given only knowledge of the organism's stoichiometry. Knowledge of reaction mechanisms can be useful for bounding fluxes, but is not necessary. 

There are a number of connections drawn in the literature between the behavior of regulatory networks and an engineered control system. Indeed, it is suggested that the only regulatory networks that should emerge in a natural biological system are those that are close to an optimal controller for maximizing some biological criterion. Hence, the expected behavior of some regulatory network should map elegantly onto the theory of optimal control, which, as a well-formulated theory, provides both descriptive ability for modelling a regulatory mechanism as well as predictive power for 1) forecasting how such mechanisms will behave in a number of environmental conditions 2) determining which regulatory mechanisms are realistic out of a number of potential candidates, given the assumption that they should be near-optimal controllers.

Another promising area that is underexplored in the existing literature is the application of stochastic hybrid systems to intra-cellular modelling (and more broadly, biological modelling in general). Similar to the benefits achieved by employing nondeterministic representation of parameters, stochastic variables are able to represent the noise inherent in biological systems 

Advantages of our modelling approach:
\begin{enumerate}
    \item Possesses the ability to construct rigorous metabolic models using only knowledge of stoichiometry
    \item Is highly generalizable, uses the canonical formalism for hybrid automata such that it can be plugged into a number of standard solvers and does not necessitate building bespoke tools
    \item Harnesses nondeterminism to allow exhaustive analysis of metabolic capacity
    \item Using the theory of extreme pathways and eigenpathways, can easily be extended to incorporate regulation without loss of generality
    \item Provides a natural formalism in which to apply control theoretic techniques to regulation, that also incorporates a metabolic component
    \item Hence, this is a metabolic-regulatory model that can be dynamically constructed from minimal information about the organism at hand (outside of stoichiometric data, which is widely available) and can be integrated with a control-theoretic component to explore regulatory mechanisms, or left nondeterministic to exhaustively explore emergent metabolic capacity
    \item Is highly visualizable, lends itself to visuals for better understanding of the model's capacity as well as identification of structural similarity with other models
\end{enumerate}



%In addition, the transformation algorithm constitutes a full \textit{reduction} from the dynamical system defined by dynamic FBA to hybrid automaton. In plain terms, any trajectory achieved by the dynamic FBA can also be achieved by the automaton. We discuss in later sections how the hybrid automaton model of the metabolism of an individual microbe can be extended to model the metabolism in a microbial co-culture.